\section{Route-A evaluation and reproducibility}

The evaluation separates exact arithmetic from the stronger dynamical and
operator claims required by Hilbert--P\'olya.

\begin{table}[t]
\centering
\small
\caption{Object-wise Route-A ruling.  ``Fail'' concerns the proposed
Hilbert--P\'olya route; ``NT'' means that the required object is not defined.}
\begin{tabular}{@{}
>{\raggedright\arraybackslash}p{0.17\textwidth}
>{\raggedright\arraybackslash}p{0.16\textwidth}
>{\raggedright\arraybackslash}p{0.18\textwidth}
>{\raggedright\arraybackslash}p{0.17\textwidth}
>{\raggedright\arraybackslash}p{0.16\textwidth}@{}}
\toprule
object & A1: primitive & A2: determinant & A3: divisor & A4: operator \\
\midrule
rational cusp & fail: zero loop clock & fail: no periodic trace & fail: shifted quotient & NT \\
full boundary & pass: classical classes & NT: period set only & fail as new RH route & NT \\
bare products & NT: no sourced dynamics & NT: no orbit determinant & fail: shifted quotient & NT \\
projected paths & NT: no primitive law & NT: no Fredholm kernel & NT: no global divisor & NT \\
\bottomrule
\end{tabular}
\label{tab:route-a}
\end{table}

For the rational-endpoint branch, the open arrows are geometric but the
section-induced loop clock is zero and the open set has no canonical
primitive-power monoid.  The full-boundary enlargement passes A1 only by
recovering classical hyperbolic conjugacy classes.  We prove that its
automorphy periods are signed Selberg lengths; we do not infer a determinant
from that period set.

For bare squarefree scattering products, commutativity does not by itself
forbid a finite or Fredholm determinant.  The A1/A2 defects are instead that
no source-derived evolution, primitive repetition law, or transfer kernel has
been defined.  Conditional on encoding steps by \(\Phi_N(s_j)\), the frozen
matrix product is permutation-invariant.  The projector-resolved witness
escapes that algebra but remains not testable at every Route-A gate beyond its
finite assignment/path-sensitivity check.

\paragraph{A3: arithmetic divisor.}
The open coefficients and scattering matrices contain exact zeta arithmetic.
Each standard eigenchannel still contains
\(\Lambda(2s-1)/\Lambda(2s)\), multiplied by local factors whose zeros and
poles lie on \(\Re s=1\) and \(\Re s=0\).  It therefore retains transported
poles at \(\rho/2\) and zeros at \((1+\rho)/2\), not one entire completed-zeta
divisor.

\paragraph{Operator boundary.}
The hyperbolic Laplacian is self-adjoint and physical-line scattering is
unitary.  Those classical facts do not turn scattering resonances into a
discrete point spectrum with ordinates of the Riemann zeros.  The project
defines no new Hilbert space, operator domain, or same-clock quantization.

\subsection{Computational audit}

The producer and checker implement three independent ledgers.

\begin{enumerate}
\item For every \(q\le2000\), exact integer arithmetic computes
\(\varphi(q)\), all roots of \(-1\), and
\(n_q=(\varphi(q)+s_q)/2\).  High-precision partial sums are compared with
\cref{eq:open-zeta} in its half-plane of convergence.
\item Exact \(\SL\) arithmetic checks the failed double-coset composition and
the primitive/affine rational-section identities and their cocycle
composition laws.
\item At levels \(2,6,30,210\), direct Kronecker matrices are compared with
Walsh reconstruction, functional equations, physical-line unitarity,
determinants, spectral-parameter product permutations, and the projected
level-six witness.  Three physical-line and two off-line points are used;
calculations use 80 decimal digits.
\end{enumerate}

The independent checker imports no producer code and recomputes the compact
results at higher precision.  The unit suite includes schema validation and
tamper rejection.  No Riemann-zero table, fitted spectral scale, or learned
transition matrix is used.

The computational statements are regression certificates, not proofs of the
universal theorems.  All universal conclusions follow from the elementary
arguments in the preceding sections.
